% =====================================================================
%  Annexe A — Dictionnaire des variables
% =====================================================================
\section{Dictionnaire des variables}
\label{annexe:variables}

Positions et unités d'après la documentation de la base~\cite{cool_welddb} ; noms
utilisés dans les notebooks. Colonne \og Usage \fg{} : E11/E13/E16 = entrée des
schémas 11, 13 ou 16 (\cref{tab:schemas}) ; C = cible ; X = non utilisée.

{\footnotesize
\begin{longtable}{@{}r l >{\raggedright\arraybackslash}p{4.7cm} >{\raggedright\arraybackslash}p{2.6cm} >{\raggedright\arraybackslash}p{2.4cm}@{}}
\caption{Les 44 champs de la Cambridge Weld Database.}\label{tab:dictionnaire}\\
\toprule
Pos. & Nom (notebooks) & Signification & Unité & Usage \\
\midrule
\endfirsthead
\toprule
Pos. & Nom (notebooks) & Signification & Unité & Usage \\
\midrule
\endhead
\bottomrule
\endfoot
1  & \texttt{C}  & Carbone   & \% massique & E11 \\
2  & \texttt{Si} & Silicium  & \% massique & E11 \\
3  & \texttt{Mn} & Manganèse & \% massique & E11 \\
4  & \texttt{S}  & Soufre    & \% massique & E11 \\
5  & \texttt{P}  & Phosphore & \% massique & E11 \\
6  & \texttt{Ni} & Nickel    & \% massique & E16 \\
7  & \texttt{Cr} & Chrome    & \% massique & E16 \\
8  & \texttt{Mo} & Molybdène & \% massique & E16 \\
9  & \texttt{V}  & Vanadium  & \% massique & X \\
10 & \texttt{Cu} & Cuivre    & \% massique & X \\
11 & \texttt{Co} & Cobalt    & \% massique & X \\
12 & \texttt{W}  & Tungstène & \% massique & X \\
13 & \texttt{O\_ppm}  & Oxygène   & ppm & E13 \\
14 & \texttt{Ti\_ppm} & Titane    & ppm (unité à revoir) & X \\
15 & \texttt{N\_ppm}  & Azote     & ppm & E13 \\
16 & \texttt{Al\_ppm} & Aluminium & ppm (unité à revoir) & X \\
17 & \texttt{B\_ppm}  & Bore      & ppm & X \\
18 & \texttt{Nb\_ppm} & Niobium   & ppm & X \\
19 & \texttt{Sn\_ppm} & Étain     & ppm (unité à revoir) & X \\
20 & \texttt{As\_ppm} & Arsenic   & ppm (unité à revoir) & X \\
21 & \texttt{Sb\_ppm} & Antimoine & ppm (unité à revoir) & X \\
22 & \texttt{current\_A} & Courant de soudage & A & E11 \\
23 & \texttt{voltage\_V} & Tension            & V & E11 \\
24 & \texttt{ac\_dc}     & Type de courant (AC/DC) & catégoriel & X \\
25 & \texttt{polarity}  & Polarité de l'électrode & catégoriel & X \\
26 & \texttt{heat\_input\_kJmm} & Apport thermique & kJ/mm & E11 \\
27 & \texttt{interpass\_C} & Température entre passes & °C & E11 \\
28 & \texttt{weld\_type} & Procédé de soudage & catégoriel & tous schémas \\
29 & \texttt{pwht\_temp\_C} & Température du traitement après soudage & °C & E11 \\
30 & \texttt{pwht\_time\_h} & Durée du traitement après soudage & h & E11 \\
31 & \texttt{yield\_strength} & Limite d'élasticité & MPa & C \\
32 & \texttt{uts} & Résistance maximale à la traction & MPa & C \\
33 & \texttt{elongation} & Allongement & \% & C \\
34 & \texttt{reduction\_area} & Réduction de section & \% & X (complémentaire) \\
35 & \texttt{charpy\_temp\_C} & Température d'essai Charpy & °C & entrée Charpy \\
36 & \texttt{charpy\_J} & Énergie Charpy & J & C (exploratoire) \\
37 & \texttt{hardness} & Dureté (protocole Vickers séparé) & kg/mm$^2$ & X \\
38 & \texttt{fatt50} & Température de transition à 50\,\% de rupture ductile & non documentée & X \\
39 & \texttt{prim\_ferrite} & Ferrite primaire & \% & X \\
40 & \texttt{ferrite\_2nd\_phase} & Ferrite à seconde phase & \% & X \\
41 & \texttt{acicular\_ferrite} & Ferrite aciculaire & \% & X \\
42 & \texttt{martensite} & Martensite & \% & X \\
43 & \texttt{ferrite\_carbide} & Agrégat ferrite-carbure & \% & X \\
44 & \texttt{weld\_id} & Identifiant de la soudure et de sa source & texte & X (sert aux groupes) \\
\end{longtable}
}

Dans la table maître, le procédé est stocké en \texttt{weld\_type\_raw} (brut)
et \texttt{weld\_process} (synonymes regroupés) ; la dureté en
\texttt{hardness\_value} et \texttt{hardness\_protocol}. Les colonnes
\texttt{n\_unavail\_A1\_13}, \texttt{in\_strict13} et \texttt{has\_<cible>}
définissent la cohorte et l'éligibilité de chaque cible.
